% p005_b 的电脑重画（① 光电效应发生电路：光电管（左板斜放、右板竖直）与电流表 A 串联成回路）
\begin{tikzpicture}[line width=0.7pt, >={Stealth[length=2.2mm]}]
  % 回路导线与电流表
  \node[draw, circle, inner sep=0pt, minimum size=0.84cm] (A) at (3.1,0) {A};
  \draw (1.39,2.7) -- (0,2.7) -- (0,0) -- (A.west);
  \draw (A.east) -- (6.2,0) -- (6.2,2.7) -- (4.35,2.7);
  % 两个极板
  \draw (1.1,3.5) -- (1.65,2.0);
  \draw (4.35,3.75) -- (4.35,1.8);
  % 入射光与电子
  \draw[->] (3.45,4.4) -- (2.6,3.95);
  \node at (4.1,4.2) {$h\nu$};
  \node at (2.05,3.45) {$e^-$};
  \draw[->] (2.85,3.1) -- (3.7,3.1);
  % 管内的回形箭头及右侧标注
  \draw[->] (1.0,1.05) .. controls (1.9,0.78) and (3.1,0.85) .. (3.9,0.85)
        .. controls (4.7,0.85) and (4.75,1.95) .. (3.9,1.95) -- (2.75,1.95);
  \node[align=center] at (5.1,0.85) {$i$\\方向};
\end{tikzpicture}
